\documentclass[12pt,a4paper]{article}

% =====================
% Packages
% =====================
\usepackage[utf8]{inputenc}
\usepackage[T1]{fontenc}
\usepackage[french]{babel}
\usepackage{amsmath,amssymb}
\usepackage{geometry}
\usepackage{enumitem}
\usepackage{xcolor}
\usepackage{multicol}

\geometry{margin=2cm}
\setlength{\parindent}{0pt}

% =====================
% Document
% =====================
\begin{document}

\begin{center}
    {\Large \textbf{Mathématiques – QCM}}\\[2mm]
    {\large Première STMG}\\[2mm]
    \textit{Épreuve anticipée – Automatismes}
\end{center}

\vspace{0.4cm}
\hrule
\vspace{0.5cm}

\textbf{Consigne :} Pour chaque question, une seule réponse est correcte.  
Aucune justification n’est demandée.

\vspace{0.5cm}

% =====================
% QCM
% =====================

\begin{enumerate}

% ---------------------
\item \textbf{Pourcentages}

\begin{enumerate}[label=\arabic{*}.\arabic*]
\item L’opération qui permet de calculer \textbf{25 \% de 480} est :
\begin{multicols}{2}
\begin{enumerate}[label=\alph*)]
\item $\dfrac{480}{25 \times 100}$
\item $25 \times 480 \times 0{,}1$
\item $\dfrac{480 \times 100}{25}$
\item $\dfrac{1}{4} \times 480$
\end{enumerate}
\end{multicols}

\item Un article coûte 120 €. Il augmente de \textbf{10 \%}, puis encore de \textbf{10 \%}.  
L’augmentation totale est :
\begin{multicols}{2}
\begin{enumerate}[label=\alph*)]
\item $(10 \%)^2$
\item $19 \%$
\item $20 \%$
\item $21 \%$
\end{enumerate}
\end{multicols}
\end{enumerate}

% ---------------------
\item \textbf{Comparaison et ordre de grandeur}

\begin{enumerate}[label=\arabic{*}.\arabic*]
\item On considère les nombres :
\[
A=\frac{1}{5}, \quad B=\frac{19}{100}, \quad C=0{,}21
\]
Le classement par ordre croissant est :
\begin{multicols}{2}
\begin{enumerate}[label=\alph*)]
\item $A < B < C$
\item $A < C < B$
\item $B < A < C$
\item $C < B < A$
\end{enumerate}
\end{multicols}

\item Parmi les nombres suivants, le \textbf{plus grand} est :
\[
A=\left(\frac{1}{5}\right)^2 \quad
B=\left(\frac{1}{2}\right)^5 \quad
C=0{,}05 \quad
D=\left(\frac{1}{3}\right)^3
\]
\begin{multicols}{2}
\begin{enumerate}[label=\alph*)]
\item $A$
\item $B$
\item $C$
\item $D$
\end{enumerate}
\end{multicols}
\end{enumerate}

% ---------------------
\item \textbf{Fractions et proportions}

\begin{enumerate}[label=\arabic{*}.\arabic*]
\item Le tiers d’un quart correspond à la fraction :
\begin{multicols}{2}
\begin{enumerate}[label=\alph*)]
\item $\dfrac{1}{7}$
\item $\dfrac{3}{4}$
\item $\dfrac{1}{12}$
\item $1$
\end{enumerate}
\end{multicols}
\end{enumerate}

% ---------------------
\item \textbf{Calcul numérique}

\begin{enumerate}[label=\arabic{*}.\arabic*]
\item On considère :
\[
A = 10 + 0{,}1 + \frac{1}{1000}
\]
Alors :
\begin{multicols}{2}
\begin{enumerate}[label=\alph*)]
\item $A = \dfrac{20 - 1}{1000}$
\item $A = 10{,}101$
\item $A = 10{,}110$
\item $A = \dfrac{10{,}110}{1000}$
\end{enumerate}
\end{multicols}
\end{enumerate}

% ---------------------
\item \textbf{Fractions et décimaux}

\begin{enumerate}[label=\arabic{*}.\arabic*]
\item Quelle expression est équivalente à \textbf{0,21} ?
\begin{multicols}{2}
\begin{enumerate}[label=\alph*)]
\item $\dfrac{21}{100}$
\item $\dfrac{21}{10}$
\item $\dfrac{210}{1000}$
\item $2{,}1$
\end{enumerate}
\end{multicols}
\end{enumerate}

% ---------------------
\item \textbf{Évolution de quantités}

\begin{enumerate}[label=\arabic{*}.\arabic*]
\item Une quantité diminue de \textbf{10 \%}, puis encore de \textbf{10 \%}.  
L’évolution globale est :
\begin{multicols}{2}
\begin{enumerate}[label=\alph*)]
\item une baisse de $20 \%$
\item une baisse inférieure à $20 \%$
\item une baisse exactement égale à $20 \%$
\item une baisse supérieure à $20 \%$
\end{enumerate}
\end{multicols}
\end{enumerate}

% ---------------------
\item \textbf{Priorités de calcul}

\begin{enumerate}[label=\arabic{*}.\arabic*]
\item Quelle écriture correspond à :
\[
3 + 5 \times 2
\]
\begin{multicols}{2}
\begin{enumerate}[label=\alph*)]
\item $(3 + 5) \times 2$
\item $3 + 5 \times 2$
\item $(3 + 5 \times 2) / 2$
\item $3 \times 5 + 2$
\end{enumerate}
\end{multicols}
\end{enumerate}

% ---------------------
\item \textbf{Probabilités}

\begin{enumerate}[label=\arabic{*}.\arabic*]
\item Un événement a une probabilité égale à $\dfrac{1}{4}$.  
La probabilité de l’événement contraire est :
\begin{multicols}{2}
\begin{enumerate}[label=\alph*)]
\item $\dfrac{1}{4}$
\item $\dfrac{3}{4}$
\item $4$
\item $\dfrac{5}{4}$
\end{enumerate}
\end{multicols}
\end{enumerate}

\end{enumerate}

\vspace{0.5cm}
\textbf{Réponses :} \dotfill

\end{document}
